% Post-CP13 research source; deliberately unassembled during the CP14 freeze.
\begin{example}[Refuted: ample Chow sections preserve codimension-two boundaries]
\label{D:Chow-surface-boundary-obstruction}
Let $k$ be any algebraically closed field, let
$S=\mathbf P^3_k$, and let $T=L=V(x_0,x_1)$ be a line.
There is a smooth projective modification $\pi:Y\to S$ and a
very ample invertible sheaf $\mathcal H$ on $Y$ such that the
normal image of a geometric generic hyperplane section has a
divisorial intersection with $L$.  Indeed, this happens for
every integral effective ample divisor on $Y$ other than the
exceptional divisor.  Thus an arbitrary ample-section argument
on a projective Chow modification does not preserve the
codimension-two boundary condition required in D.2(ii).
\end{example}

\begin{proof}
Take $Y=\operatorname{Bl}_L\mathbf P^3$.  The affine blowup
charts identify it with the incidence hypersurface
\[
 Y=V(x_0v-x_1u)\subset\mathbf P^3\times\mathbf P^1,
 \qquad [u:v]\in\mathbf P^1.
\]
For example, on $u\ne0$ the equation is $x_1=x_0(v/u)$,
which is the chart of the blowup in the generator $x_0$;
the other chart uses $x_1$.  These charts also show smoothness
\cite[Tag 0804, Lemma 31.33.2]{StacksDChowSection}.
The exceptional divisor is
$E=\pi^{-1}(L)=L\times\mathbf P^1$.
The restriction
$\mathcal H=\mathcal O_{\mathbf P^3\times\mathbf P^1}(1,1)|_Y$
is very ample by the explicit Segre embedding
\cite[Tag 01WD, Lemma 27.13.6]{StacksDChowSection}.
Writing $H=\pi^*\mathcal O_{\mathbf P^3}(1)$, one also has
$\mathcal H=2H-E$.

Every section of $\mathcal H$ has the form
$u\,a(x)+v\,b(x)$ for linear forms $a,b$, modulo the incidence
equation.  To check completeness, use
\[
 0\longrightarrow\mathcal O
 \longrightarrow\mathcal O(1,1)
 \longrightarrow\mathcal H\longrightarrow0
\]
on $\mathbf P^3\times\mathbf P^1$ and the vanishing of
$H^1(\mathcal O)$, obtained from the standard projective-space
cohomology calculation
\cite[Tag 01XT, Lemma 30.8.1]{StacksDChowSection}.
Off $L$, its image is the quadric
\[
                 x_0a(x)+x_1b(x)=0,
\]
which contains $L$.
The family contains the smooth quadric
\[
 Q=V(x_0x_2+x_1x_3),\qquad
 \widetilde Q=Y\cap V(ux_2+vx_3).
\]
Smoothness holds in every characteristic, since the four
partial derivatives are $x_2,x_3,x_0,x_1$.  The projection
$\widetilde Q\to Q$ is an isomorphism: away from $L$ its
inverse has $[u:v]=[x_0:x_1]$, and near $L$ it has
$[u:v]=[-x_3:x_2]$.  These formulas agree on the overlap by
the equation of $Q$.  In particular,
$\widetilde Q\cap E\to L$ is an isomorphism.
The smooth locus of the projective family of quadrics is
open; the displayed smooth member shows that its geometric
generic member is smooth.  Its line $L$ is Cartier, and its
strict transform is likewise isomorphic to it.  Therefore
normalizing the image does not remove the boundary curve.

There is a numerical explanation independent of this chosen
linear system.  Let $B\subset Y$ be an integral effective
ample divisor different from $E$.  On every exceptional
fibre $F\simeq\mathbf P^1$, the restriction of
$\mathcal O_Y(B)$ has positive degree.  Its defining section
either vanishes identically on $F$ or has a zero there.
Consequently $B$ meets every such fibre, and
$L\subset\pi(B)$.  Since $B\ne E$, its image is a surface.
The finite normalization of that image still has a
one-dimensional inverse image of $L$.  This proves the
assertion for every such ample divisor and, in particular,
for geometric generic hyperplanes of any projective
embedding of $Y$.
\end{proof}

\begin{remark}[Scope and a possible repair]
This example refutes the stated slicing inference, not D.1,
D.2(ii), or the existence of every possible surface reduction.
Indeed, $S$ itself is projective.  A general divisor from the
non-ample system $|H|$ maps to a plane meeting $L$ in one
point, and therefore repairs the boundary dimension in this
example.  The difficulty over a general proper nonprojective
base is obtaining a sufficiently mobile system with the
required zero degree on modification fibres.  Moreover,
connectedness of the restricted original covers alone would
not suffice for D.1: their ramification along divisors newly
created over a codimension-two boundary has no supplied
tameness condition.

The exact symbolic and finite-field certificate is
\nolinkurl{parts/D/code/check_Chow_surface_boundary.py}; its
output is
\nolinkurl{parts/D/output/Chow_surface_boundary.json}.
It verifies the elimination identities, the seven-dimensional
quadrics-through-$L$ linear system, smoothness equations,
and the stated incidence fibres over $\mathbf F_2$,
$\mathbf F_3$, and $\mathbf F_5$.
\end{remark}
